\section{Data Analysis III: What to Teach First (RQ4)}\label{sec:rq4}

\subsection{Model}

RQ4 turns the evidence into a decision. For each candidate course $c$ the model chooses whether to
run it ($r_c \in \{0,1\}$) and how many seats to offer ($s_c$). For each skill area $k$ it counts
learners whose gap is closed, $z_k$:
\begin{align*}
\max \; & \textstyle\sum_k w_k z_k \\
\text{s.t. } & z_k \le \min\!\big(G_k,\ \textstyle\sum_c a_{ck}\, s_c\big) && \text{no value beyond the gap}\\
& m_c r_c \le s_c \le M_c r_c && \text{batch size and seat limits}\\
& \textstyle\sum_c (F_c r_c + p_c s_c) \le B && \text{budget}\\
& \textstyle\sum_c h_c r_c \le H && \text{trainer-hours}\\
& r_c \le r_{c'} \ \text{for prerequisites } c' \text{ of } c; \quad \textstyle\sum_{c\in g} r_c \le 1 && \text{either/or groups}
\end{align*}
$G_k$ (gap) and $w_k$ (value) come from the data; course parameters $F_c, p_c, h_c, M_c, a_{ck}$ are
documented assumptions. The model is solved exactly with CP-SAT \citep{cpsat}.

\begin{table}[H]\centering\small
\caption{RQ4 inputs: what is data and what is assumed.}\label{tab:rq4in}
\begin{tabularx}{\linewidth}{@{}L{3.1cm}L{4.6cm}Y@{}}
\toprule
Input & Source & Values \\
\midrule
Gap $G_k$ & \textbf{Data (JDS):} juniors below the high-hike group's median & Maths 82 \textbullet{} Coding 77 \textbullet{} AI/ML 71 \textbullet{} Big data 65 \textbullet{} Storytelling 58\\
Value $w_k$ & \textbf{Data (RQ1+RQ2):} equal mix of demand share, log market OR, log RQ2 OR, each scaled to its maximum & Maths 100 \textbullet{} Coding 99.7 \textbullet{} AI/ML 77.3 \textbullet{} Storytelling 62.9 \textbullet{} Big data 59.4\\
Uncertainty & \textbf{Data:} CI lower/upper bounds $\times$ 3 weighting schemes & 9 scenarios\\
Courses & \textbf{Assumption}, checked against public sources & 9 courses (Appendix~\ref{app:catalogue})\\
Budget, capacity & \textbf{Assumption} (scenario) & Rs~4 lakh, 200 trainer-hours\\
\bottomrule
\end{tabularx}
\end{table}

Cost assumptions: trainers at Rs~1,430--2,000 per hour, the upper-mid of the Rs~250--2,000 range
reported for Indian data-science trainers \citep{knowledgehut} (deliberately conservative); SAS
Viya for Learners \citep{sasvfl}, Tableau for students and Power BI Desktop are free, so per-seat
costs are materials only; GPU time for deep learning at about Rs~100 per hour \citep{gpuindia2026}.
Personality (RQ3) is not modelled as a course, because traits are relatively stable; it informs
mentoring and assessment instead (Section~\ref{sec:implications}).

\subsection{The optimal plan}

\begin{table}[H]\centering\small
\caption{Proven-optimal plan for Rs~4 lakh and 200 trainer-hours.}\label{tab:plan}
\begin{tabular}{@{}l r r r@{}}
\toprule
Course & Seats & Cost (Rs lakh) & Trainer-hours\\
\midrule
Applied Statistics \& Probability & 60 & 0.78 & 40\\
SAS Programming on SAS Viya for Learners & 60 & 0.45 & 30\\
Machine Learning Foundations & 50 & 0.97 & 48\\
Python \& SQL for Analytics & 35 & 0.71 & 40\\
Data Storytelling with Tableau & 58 & 0.62 & 30\\
\midrule
Total & 263 & 3.53 & 188\\
\bottomrule
\end{tabular}
\end{table}

\fig{fig_rq4_plan.png}{Optimal plan (left) and the share of each skill gap it closes (right).}

The plan closes 100\% of the coding, maths \& stats and storytelling gaps and 70\% of the AI \& ML
gap (ML Foundations is capped at 50 seats). It leaves the big-data gap open.

\textbf{Why other courses were left out.} Forcing each excluded course into the plan and re-solving
shows its cost: Integrated Bootcamp $-20.7\%$ impact (high fixed cost, 90 trainer-hours, partial
coverage), Deep Learning \& NLP $-12.8\%$ (needs ML first and 40 more hours), Big-Data Stack $-5.7\%$
(lowest-valued area, does not fit the remaining hours), Power BI $-1.1\%$ (either/or with Tableau,
which closes the full gap at the same per-seat cost).

\subsection{Sensitivity: what would change the decision?}

\textbf{Money is not the constraint}: Rs~0.47 lakh remains unspent, and an extra Rs~1 lakh adds no
impact. \textbf{Trainer capacity is}, but in whole-course steps. The linear relaxation suggests each
extra trainer-hour is worth 27.5 impact units, yet the exact integer re-solve with +40 hours adds
nothing, because the next course (Big Data) needs 40 hours \emph{and} Rs~1.1 lakh at the same time.
This is precisely why integer optimisation, not a linear approximation, is needed. With a generous
budget, extra trainer capacity buys:

\begin{table}[H]\centering\small
\caption{What additional trainer-hours buy (budget not limiting).}\label{tab:hours}
\begin{tabular}{@{}r l r l@{}}
\toprule
Trainer-hours & Added to the plan & Impact (change) & Gap closed\\
\midrule
200 & base plan (5 courses) & 23,391 & AI/ML 70\%, big data 0\%\\
240 & + Big-Data stack & 25,767 (+10.2\%) & big data 62\%\\
280 & + Deep Learning \& NLP & 27,158 (+16.1\%) & AI/ML 96\%\\
400 & + Bootcamp & 27,389 (+17.1\%) & gains level off\\
\bottomrule
\end{tabular}
\end{table}

\fig[0.8\linewidth]{fig_rq4_pareto.png}{Cost--impact Pareto frontier with 200 and 400
trainer-hours.}

\subsection{Robustness}

The plan was re-optimised under nine evidence scenarios: balanced, market-led and outcome-led
weightings, each at the lower bound, point estimate and upper bound of every confidence interval
(Figure~\ref{fig:fig_rq4_robust.png}). \textbf{Statistics (60), SAS (60), ML Foundations (50) and
Python \& SQL (35) are chosen with identical seat counts in all nine scenarios}, forming a safe core.
Tableau storytelling is chosen in 7 of 9 and is replaced by the Big-Data stack only when the
evidence is weighted purely by what job postings advertise.

\fig{fig_rq4_robust.png}{Seats per course across nine evidence scenarios.}
